% GENERATED FILE. Edit canonical lessons, then run npm run generate:books.
% canonical-source-hash: fnv1a64:ec129be7998f4202
% canonical-lessons: AR-C145-malih, AR-C145-fakhm, AR-C145-basit, AR-C145-murih, AR-C145-mashhur

\chapter{Salty, Luxurious, Simple, Comfortable, Famous}
\label{ch:ar-salty-luxurious-simple-comfortable-famous}
\hlchaptermodality{145}

\begin{chapteropening}
\textbf{By the end of this chapter:} \emph{I can say salty, luxurious, simple, comfortable and famous.}

Complete the last lesson of chapter 145: I can say salty, luxurious, simple, comfortable and famous.
\end{chapteropening}

\section[\texorpdfstring{\ar{مالح} (māliḥ)}{مالح (māliḥ)}]{\ar{مالح} (māliḥ) — salty}

\label{lesson:AR-C145-malih}



\begin{quote}
Before the new one: say the Arabic for a farmer, then the Arabic for a university.
\end{quote}

\subsection*{You'll want to know: \ar{مالح}}
\textbf{\ar{مالح}} — \emph{māliḥ} — \textquotedblleft{}salty\textquotedblright{}.

\begin{cousinweb}
\textbf{\ar{مالح}} (\emph{māliḥ}) is built on the native Arabic root \emph{m-l-ḥ}, \textquotedblleft{}salt\textquotedblright{} (Lane). \emph{Māliḥ} is \textquotedblleft{}salty\textquotedblright{}.
\end{cousinweb}

\begin{grammarlens}[title={What you've built}]
One more word for this chapter.
\end{grammarlens}

\subsection*{Your turn}
\begin{itemize}
\raggedright
  \item \emph{Say it:} \emph{māliḥ}
  \item \emph{Say it:} \emph{māliḥ}, once more
  \item \emph{Say it:} say \emph{māliḥ}
  \item \emph{Recall:} say the Arabic for a farmer, then the Arabic for a university, then say \emph{māliḥ} again
\end{itemize}

\begin{tcolorbox}[breakable,colback=teal!4,colframe=teal!35!black,title={Before you move on}]
What does \ar{مالح} mean? (\textquotedblleft{}Salty\textquotedblright{}.) Say it once more.
\end{tcolorbox}
\section[\texorpdfstring{\ar{فخم} (fakhm)}{فخم (fakhm)}]{\ar{فخم} (fakhm) — luxurious}

\label{lesson:AR-C145-fakhm}



\begin{quote}
Before the new one: say the Arabic for a university, then the Arabic for salty.
\end{quote}

\subsection*{You'll want to know: \ar{فخم}}
\textbf{\ar{فخم}} — \emph{fakhm} — \textquotedblleft{}luxurious\textquotedblright{}.

\begin{cousinweb}
\textbf{\ar{فخم}} (\emph{fakhm}) is built on the native Arabic root \emph{f-kh-m}, \textquotedblleft{}to be large, to be grand\textquotedblright{} (Lane). \emph{Fakhm} is \textquotedblleft{}grand\textquotedblright{}, hence \textquotedblleft{}luxurious\textquotedblright{}.
\end{cousinweb}

\begin{grammarlens}[title={What you've built}]
One more word for this chapter.
\end{grammarlens}

\subsection*{Your turn}
\begin{itemize}
\raggedright
  \item \emph{Say it:} \emph{fakhm}
  \item \emph{Say it:} \emph{fakhm}, once more
  \item \emph{Say it:} say \emph{fakhm}
  \item \emph{Recall:} say the Arabic for a university, then the Arabic for salty, then say \emph{fakhm} again
\end{itemize}

\begin{tcolorbox}[breakable,colback=teal!4,colframe=teal!35!black,title={Before you move on}]
What does \ar{فخم} mean? (\textquotedblleft{}Luxurious\textquotedblright{}.) Say it once more.
\end{tcolorbox}
\section[\texorpdfstring{\ar{بسيط} (basīṭ)}{بسيط (basīṭ)}]{\ar{بسيط} (basīṭ) — simple}

\label{lesson:AR-C145-basit}



\begin{quote}
Before the new one: say the Arabic for salty, then the Arabic for luxurious.
\end{quote}

\subsection*{You'll want to know: \ar{بسيط}}
\textbf{\ar{بسيط}} — \emph{basīṭ} — \textquotedblleft{}simple\textquotedblright{}.

\begin{cousinweb}
\textbf{\ar{بسيط}} (\emph{basīṭ}) is built on the native Arabic root \emph{b-s-ṭ}, \textquotedblleft{}to spread out\textquotedblright{} (Lane). \emph{Basīṭ} is \textquotedblleft{}spread out, plain\textquotedblright{}, hence \textquotedblleft{}simple\textquotedblright{}.
\end{cousinweb}

\begin{grammarlens}[title={What you've built}]
One more word for this chapter.
\end{grammarlens}

\subsection*{Your turn}
\begin{itemize}
\raggedright
  \item \emph{Say it:} \emph{basīṭ}
  \item \emph{Say it:} \emph{basīṭ}, once more
  \item \emph{Say it:} say \emph{basīṭ}
  \item \emph{Recall:} say the Arabic for salty, then the Arabic for luxurious, then say \emph{basīṭ} again
\end{itemize}

\begin{tcolorbox}[breakable,colback=teal!4,colframe=teal!35!black,title={Before you move on}]
What does \ar{بسيط} mean? (\textquotedblleft{}Simple\textquotedblright{}.) Say it once more.
\end{tcolorbox}
\section[\texorpdfstring{\ar{مريح} (murīḥ)}{مريح (murīḥ)}]{\ar{مريح} (murīḥ) — comfortable}

\label{lesson:AR-C145-murih}



\begin{quote}
Before the new one: say the Arabic for luxurious, then the Arabic for simple.
\end{quote}

\subsection*{You'll want to know: \ar{مريح}}
\textbf{\ar{مريح}} — \emph{murīḥ} — \textquotedblleft{}comfortable\textquotedblright{}.

\begin{cousinweb}
\textbf{\ar{مريح}} (\emph{murīḥ}) is built on the native Arabic root \emph{r-w-ḥ}, \textquotedblleft{}rest; wind\textquotedblright{} (Lane). \emph{Murīḥ} is the Form IV participle, \textquotedblleft{}giving rest\textquotedblright{}.
\end{cousinweb}

\begin{grammarlens}[title={What you've built}]
One more word for this chapter.
\end{grammarlens}

\subsection*{Your turn}
\begin{itemize}
\raggedright
  \item \emph{Say it:} \emph{murīḥ}
  \item \emph{Say it:} \emph{murīḥ}, once more
  \item \emph{Say it:} say \emph{murīḥ}
  \item \emph{Recall:} say the Arabic for luxurious, then the Arabic for simple, then say \emph{murīḥ} again
\end{itemize}

\begin{tcolorbox}[breakable,colback=teal!4,colframe=teal!35!black,title={Before you move on}]
What does \ar{مريح} mean? (\textquotedblleft{}Comfortable\textquotedblright{}.) Say it once more.
\end{tcolorbox}
\section[\texorpdfstring{\ar{مشهور} (mashhūr)}{مشهور (mashhūr)}]{\ar{مشهور} (mashhūr) — famous}

\label{lesson:AR-C145-mashhur}



\begin{quote}
Before the new one: say the Arabic for simple, then the Arabic for comfortable.
\end{quote}

\subsection*{You'll want to know: \ar{مشهور}}
\textbf{\ar{مشهور}} — \emph{mashhūr} — \textquotedblleft{}famous\textquotedblright{}.

\begin{cousinweb}
\textbf{\ar{مشهور}} (\emph{mashhūr}) is built on the native Arabic root \emph{sh-h-r}, \textquotedblleft{}to make known\textquotedblright{} (Lane). \emph{Mashhūr} is the passive participle, \textquotedblleft{}made known\textquotedblright{}, hence \textquotedblleft{}famous\textquotedblright{}.
\end{cousinweb}

\begin{grammarlens}[title={What you've built}]
One more word for this chapter.
\end{grammarlens}

\subsection*{Your turn}
\begin{itemize}
\raggedright
  \item \emph{Say it:} \emph{mashhūr}
  \item \emph{Say it:} \emph{mashhūr}, once more
  \item \emph{Say it:} say \emph{mashhūr}
  \item \emph{Recall:} say the Arabic for simple, then the Arabic for comfortable, then say \emph{mashhūr} again
\end{itemize}

\begin{tcolorbox}[breakable,colback=teal!4,colframe=teal!35!black,title={Before you move on}]
What does \ar{مشهور} mean? (\textquotedblleft{}Famous\textquotedblright{}.) Say it once more.
\end{tcolorbox}
